
\documentclass[11pt,a4paper]{article}

\usepackage[ngerman]{babel}
\usepackage[T1]{fontenc}
\usepackage[utf8]{inputenc}
\usepackage{lmodern}
\usepackage{microtype}
\usepackage[a4paper,margin=2.5cm]{geometry}
\usepackage{amsmath,amssymb,amsthm,mathtools}
\usepackage{booktabs}
\usepackage{enumitem}
\usepackage{hyperref}
\usepackage{csquotes}

\hypersetup{
  colorlinks=true,
  linkcolor=black,
  urlcolor=blue,
  pdftitle={Kanonische Mehrlingspolynome und gerichtete Uebergangsoperatoren mit Literaturapparat},
  pdfauthor={Thomas Hoffbauer}
}

\theoremstyle{definition}
\newtheorem{satz}{Satz}[section]
\newtheorem{lemma}[satz]{Lemma}
\newtheorem{korollar}[satz]{Korollar}
\newtheorem{definition}[satz]{Definition}
\newtheorem{bemerkung}[satz]{Bemerkung}
\newtheorem{vermutung}[satz]{Vermutung}
\newtheorem{arbeitsvermutung}[satz]{Arbeitsvermutung}

\title{Kanonische Mehrlingspolynome und gerichtete Uebergangsoperatoren\\
\large Endredigierte Kurzfassung mit Literaturapparat}
\author{Thomas Hoffbauer}
\date{\today}

\begin{document}
\maketitle

\begin{abstract}
Wir untersuchen zentrierte kanonische Polynome lokaler Mehrlingsmuster und definieren daraus gerichtete Uebergangsoperatoren zwischen Mehrlingsklassen. Der Schwerpunkt liegt auf Symmetrieerhalt, Diskriminantenwachstum, ungeraden Momenten und Faktorisierungsdefekten. Methodisch verbindet der Text elementare Zahlentheorie, Matrixanalyse und spektrale Graphentheorie. Numerisch zeigt sich keine robuste exklusive Naehe zu Zeta-Referenzdaten, wohl aber eine klare interne Ordnungsstruktur: symmetrieerhaltende Uebergaenge sind signifikant guenstiger als symmetriebrechende.
\end{abstract}

\section{Einleitung}

Die vorliegende Kurzfassung bewegt sich an einer Schnittstelle von elementarer Zahlentheorie, Matrixanalyse und spektraler Graphentheorie. Auf der zahlentheoretischen Seite steht die Idee, lokale Mehrlingsmuster nach Zentrierung in algebraische Normalformen zu uebersetzen; auf der operatorischen Seite werden aus diesen Mustern gerichtete Uebergangsgraphen und daraus abgeleitete Operatoren aufgebaut. Fuer klassische Grundlagen der Zahlentheorie ist Hardy--Wright weiterhin ein natuerlicher Referenzpunkt, waehrend fuer die spektrale Organisation endlicher gewichteter Graphen Chung und fuer Matrix- und Begleitmatrix-Fragen Horn--Johnson einschlaegig sind \cite{hardywright2008,chung1997,hornjohnson2012}.

Die Leitfrage lautet nicht, ob bereits eine abgeschlossene Theorie vorliegt, sondern ob sich innerhalb dieses polynomialen Strangs ein stabiler mathematischer Kern isolieren laesst. Die bisher staerkste Beobachtung ist intern: Symmetrieerhalt ordnet den Uebergangsraum deutlich.

\section{Kanonische Mehrlingspolynome}

\begin{definition}
Sei \(\omega=(a_1,\dots,a_k)\) ein geordnetes Muster. Mit
\[
m(\omega):=\frac{1}{k}\sum_{\nu=1}^k a_\nu,\qquad b_\nu:=a_\nu-m(\omega)
\]
definieren wir das zentrierte kanonische Polynom
\[
Q_\omega(y):=\prod_{\nu=1}^k (y-b_\nu).
\]
\end{definition}

\begin{lemma}
Im zentrierten Polynom \(Q_\omega(y)\) verschwindet der Koeffizient von \(y^{k-1}\).
\end{lemma}

\begin{proof}
Es gilt \(\sum_{\nu=1}^k b_\nu=0\). Der Koeffizient von \(y^{k-1}\) ist bis auf Vorzeichen genau diese Summe.
\end{proof}

\begin{lemma}
Ist die Nullstellenmenge \(\{b_1,\dots,b_k\}\) invariant unter \(b\mapsto -b\), so ist \(Q_\omega\) bei geradem Grad gerade. Bei ungeradem Grad und zusaetzlicher Nullstelle \(0\) ist \(Q_\omega\) ungerade.
\end{lemma}

\section{Uebergangskosten}

\begin{definition}
Der Symmetrieindex eines Musters \(\omega\) sei
\[
\mathfrak S(\omega):=
\begin{cases}
1,& \text{falls } Q_\omega(-y)=Q_\omega(y)\text{ oder }Q_\omega(-y)=-Q_\omega(y),\\
0,& \text{sonst.}
\end{cases}
\]
\end{definition}

\begin{definition}
Fuer einen Uebergang \(\omega_k\rightsquigarrow\omega_{k+1}\) definieren wir formal
\[
\mathcal T(\omega_k,\omega_{k+1})
=
\alpha(\mathfrak S(\omega_k)-\mathfrak S(\omega_{k+1}))^2
+\beta\,\mathfrak M_{\mathrm{odd}}(\omega_{k+1})
+\gamma\,\mathfrak D(\omega_k,\omega_{k+1})
+\delta\,(\mathfrak F(\omega_{k+1})-\mathfrak F(\omega_k)),
\]
wobei \(\mathfrak M_{\mathrm{odd}}\) den ungeraden Momentendefekt, \(\mathfrak D\) die Diskriminantenkosten und \(\mathfrak F\) einen Faktorisierungsdefekt bezeichnet.
\end{definition}

\section{Beispiele}

\begin{table}[h]
\centering
\begin{tabular}{llll}
\toprule
Uebergang & Symmetrie & Zielpolynomtyp & qualitative Kosten \\
\midrule
\((0,2)\to(0,2,6)\) & \(1\to 0\) & asymmetrischer Drilling & hoch \\
\((0,2,6,8)\to(0,2,6,8,12)\) & \(1\to 0\) & asymmetrischer Fuenfling & sehr hoch \\
\((0,2,6,8)\to(0,4,6,8,12)\) & \(1\to 1\) & symmetrischer Fuenfling & niedrig \\
\bottomrule
\end{tabular}
\caption{Charakteristische Uebergaenge zwischen Mehrlingsmustern.}
\end{table}

\section{Gerichteter Uebergangsoperator}

Wir betrachten einen gerichteten Graphen, dessen Knoten kanonische Mehrlingspolynome und dessen Kanten elementare Erweiterungen sind. Die Kantengewichte werden als exponentielle Daempfung der Uebergangskosten modelliert. Die Motivation fuer diesen Schritt ist stark von der spektralen Graphensicht inspiriert, wobei die hier verwendeten Operatoren jedoch gerichtet und daher nichtsymmetrisch sind \cite{chung1997}. Fuer die Interpretation der entstehenden Spektren und Singulaerwerte ist der matrixanalytische Rahmen natuerlich \cite{hornjohnson2012}.

\begin{arbeitsvermutung}
Im gerichteten Uebergangsgraphen kanonischer Mehrlingspolynome minimieren symmetrieerhaltende Uebergaenge im Mittel die algebraischen Uebergangskosten.
\end{arbeitsvermutung}

\section{Numerischer Befund}

Ein Vergleich der Spektren des gerichteten Operators mit \texttt{zeros6.npy} liefert keine robuste exklusive Zeta-Signatur. Dagegen zeigt die innere Statistik des Operators ein klares Muster:

\begin{table}[h]
\centering
\begin{tabular}{ccccc}
\toprule
\(\mathfrak S_{\mathrm{from}}\) & \(\mathfrak S_{\mathrm{to}}\) & Anzahl & mittlere Kosten & mittleres Gewicht \\
\midrule
0 & 0 & 1102 & 107.565586 & 0.134436 \\
0 & 1 & 130  & 5.673910   & 0.159093 \\
1 & 0 & 206  & 221.098242 & 0.003657 \\
1 & 1 & 20   & 2.554875   & 0.450206 \\
\bottomrule
\end{tabular}
\caption{Interne Statistik des Uebergangsoperators.}
\end{table}

Die externe Zeta-Referenz bleibt also derzeit methodisch schwach, waehrend die interne Symmetrieordnung stark ausfaellt. Das Resultat ist weniger als Endpunkt denn als Selektion eines belastbaren Kerns zu lesen.

\section{Diskussion und Ausblick}

Der staerkste Ertrag dieses Strangs ist gegenwaertig nicht eine aeussere spektroskopische Uebereinstimmung, sondern eine interne relationale Geometrie. Der Uebergangsoperator scheint nicht primaer Zeta-Nullstellen zu reproduzieren, sondern einen Raum guenstiger und unguenstiger Uebergaenge zwischen lokalen Polynomtypen zu ordnen. Genau darin koennte der mathematisch tragfaehige Kern liegen.

Fuer eine naechste Ausbaustufe bieten sich drei Richtungen an:
\begin{enumerate}[leftmargin=2em]
\item strengere Definition der Defektgroessen \(\mathfrak M_{\mathrm{odd}},\mathfrak D,\mathfrak F\),
\item Einbau expliziter Reduktions- und Familieninformation in den Operator,
\item und systematischer Vergleich symmetrischer und asymmetrischer Muster hoeherer Grade.
\end{enumerate}

\begin{thebibliography}{99}

\bibitem{hardywright2008}
G.~H. Hardy und E.~M. Wright.
\newblock \emph{An Introduction to the Theory of Numbers}.
\newblock 6. Auflage, Oxford University Press, 2008.

\bibitem{chung1997}
Fan R.~K. Chung.
\newblock \emph{Spectral Graph Theory}.
\newblock CBMS Regional Conference Series in Mathematics, Band~92,
American Mathematical Society, 1997.

\bibitem{hornjohnson2012}
Roger A. Horn und Charles R. Johnson.
\newblock \emph{Matrix Analysis}.
\newblock 2. Auflage, Cambridge University Press, 2012.

\end{thebibliography}

\end{document}
